\chapter{Índice de notación}
La página indicada remite al comienzo del capítulo donde se introduce o se utiliza sistemáticamente el símbolo. Las letras $X,Y$ suelen designar espacios métricos; $V,W$, espacios normados; $H$, un espacio de Hilbert. Las letras locales se definen nuevamente cuando cambia el contexto.
\small
\begin{longtable}{@{}p{0.22\textwidth}p{0.66\textwidth}r@{}}
\toprule
Símbolo & Significado & Pág.\\
\midrule
\endfirsthead
\toprule Símbolo & Significado & Pág.\\\midrule
\endhead
\bottomrule\endfoot
$\N,\Z,\Q,\R,\C$ & Naturales positivos, enteros, racionales, reales y complejos & \pageref{chap:01}\\
$\F$ & Cuerpo de escalares: $\R$ o $\C$ & \pageref{chap:01}\\
$d(x,y)$ & Distancia entre $x$ e $y$ & \pageref{chap:01}\\
$\norm{x}$ & Norma del vector $x$ & \pageref{chap:01}\\
$\norm{x}_p$ & Norma de tipo $p$ en un espacio de coordenadas o sucesiones & \pageref{chap:01}\\
$\ell^p,\ \ell^\infty$ & Espacios de sucesiones sumables a la potencia $p$ o acotadas & \pageref{chap:01}\\
$B(x,r)$ & Bola abierta de centro $x$ y radio $r>0$ & \pageref{chap:02}\\
$\overline B(x,r)$ & Bola cerrada: puntos a distancia a lo sumo $r$ de $x$ & \pageref{chap:02}\\
$d(x,A),\ \dist(x,A)$ & Ínfimo de las distancias desde $x$ a puntos de $A\ne\varnothing$ & \pageref{chap:03}\\
$\diam A$ & Diámetro de $A$: supremo de las distancias entre sus puntos & \pageref{chap:02}\\
$\interior A$ & Interior de $A$ en el espacio ambiente especificado & \pageref{chap:02}\\
$\cl A$ & Clausura de $A$; si hace falta se indica el espacio ambiente & \pageref{chap:02}\\
$\partial A$ & Frontera de $A$ & \pageref{chap:05}\\
$A'$ & Conjunto de puntos de acumulación de $A$ & \pageref{chap:03}\\
$x_n\to x$ & Convergencia de la sucesión $(x_n)$ a $x$ & \pageref{chap:03}\\
$f^{-1}(A)$ & Preimagen de $A$; no presupone que $f$ sea invertible & \pageref{chap:05}\\
$\supp f$ & Soporte: clausura del conjunto donde $f$ no se anula & \pageref{chap:05}\\
$c_{00}$ & Sucesiones con un número finito de términos no nulos & \pageref{chap:06}\\
$c,\ c_0$ & Sucesiones convergentes y sucesiones que convergen a cero & \pageref{chap:08}\\
$|r|_p$ & Valor absoluto $p$-ádico de un racional, con $p$ primo & \pageref{chap:08}\\
$\Q_p$ & Cuerpo de números $p$-ádicos, completación de $\Q$ & \pageref{chap:09}\\
$\widehat X$ & Completación de un espacio métrico $X$ & \pageref{chap:09}\\
$F_\sigma,\ G_\delta$ & Unión numerable de cerrados, intersección numerable de abiertos & \pageref{chap:10}\\
$f^n$ & Iterada de una aplicación en el contexto de puntos fijos & \pageref{chap:11}\\
$\mathbf1_A$ & Función característica de $A$, igual a $1$ en $A$ y a $0$ fuera & \pageref{chap:10}\\
$C$ & Conjunto de Cantor, en el capítulo de compactificaciones & \pageref{chap:16}\\
$[0,1]^{\N}$ & Cubo de Hilbert con la topología producto & \pageref{chap:16}\\
$\mathcal N=\N^{\N}$ & Espacio de Baire & \pageref{chap:16}\\
$A+B,\ \lambda A$ & Suma de Minkowski y dilatación de conjuntos & \pageref{chap:19}\\
$\spanop A$ & Subespacio generado por combinaciones lineales finitas de $A$ & \pageref{chap:19}\\
$\Bcal(V,W)$ & Espacio de operadores lineales acotados de $V$ en $W$ & \pageref{chap:20}\\
$\norm{T}$ & Norma de operador; supremo de $\norm{Tv}$ sobre $\norm v\le1$ & \pageref{chap:20}\\
$\ker T,\ T(V)$ & Núcleo e imagen del operador $T$ & \pageref{chap:20}\\
$V^*$ & Dual continuo del espacio normado $V$ & \pageref{chap:21}\\
$V^{**}$ & Dual del dual continuo de $V$ & \pageref{chap:21}\\
$\id,\ I$ & Aplicación identidad u operador identidad, según el contexto & \pageref{chap:21}\\
$\Kcal(V,W)$ & Operadores compactos de $V$ en $W$ & \pageref{chap:23}\\
$\operatorname{co}(A)$ & Envolvente convexa de $A$ & \pageref{chap:24}\\
$\operatorname{ext}(A)$ & Conjunto de puntos extremos de un convexo $A$ & \pageref{chap:24}\\
$C(X),\ C(K)$ & Funciones continuas; $K$ designa habitualmente un compacto & \pageref{chap:25}\\
$C_b(X)$ & Funciones continuas acotadas & \pageref{chap:25}\\
$C_c(X)$ & Funciones continuas de soporte compacto & \pageref{chap:25}\\
$C_0(X)$ & Funciones continuas que se anulan en el infinito & \pageref{chap:25}\\
$B(S,Y)$ & Funciones acotadas del conjunto $S$ al espacio métrico $Y$ & \pageref{chap:25}\\
$d_\infty(f,g)$ & Distancia uniforme: $\sup_{s\in S}d(f(s),g(s))$ & \pageref{chap:25}\\
$\norm f_\infty$ & Norma uniforme: $\sup_x\norm{f(x)}$; para escalares, $\sup_x|f(x)|$ & \pageref{chap:25}\\
$f_n\rightrightarrows f$ & Convergencia uniforme & \pageref{chap:26}\\
$\ip uv$ & Producto interno, lineal en $u$ y conjugado lineal en $v$ & \pageref{chap:29}\\
$u\perp v$ & Ortogonalidad: $\ip uv=0$ & \pageref{chap:29}\\
$M^\perp$ & Complemento ortogonal de un conjunto o subespacio $M$ & \pageref{chap:29}\\
$P_M$ & Proyección ortogonal sobre $M$, cuando existe; en Hilbert basta que $M$ sea cerrado & \pageref{chap:30}\\
$\{e_s\}_{s\in S}$ & Familia ortonormal; se especifica cuándo es una base & \pageref{chap:31}\\
$\ell^2(S)$ & Familias cuadrado sumables con índices en $S$ & \pageref{chap:29}\\
$\ell^2_w$ & Sucesiones cuadrado sumables con pesos positivos $w_n$ & \pageref{chap:29}\\
$\sum_{s\in S}a_s$ & Para $a_s\ge0$, supremo de las sumas sobre subconjuntos finitos de $S$ & \pageref{chap:29}\\
$P_n$ & Polinomio de Legendre de grado $n$ en el ejercicio correspondiente & \pageref{chap:31}\\
$T^*$ & Adjunto hilbertiano de un operador acotado $T$ & \pageref{chap:32}\\
$V/M,\ [x],\ q$ & Espacio cociente, clase $x+M$ y proyección canónica & \pageref{ampl:cocientes}\\
$p,\ \bar p,\ N_p$ & Seminorma, seminorma cociente y núcleo de $p$ & \pageref{ampl:nucleo-seminorma}\\
$\norm u_\beta$ & Norma uniforme ponderada por $e^{-\beta|t-t_0|}$ & \pageref{ampl:picard-equivalencia}\\
$S,\ u_n$ & Operador de Picard y aproximaciones sucesivas & \pageref{ampl:picard-teorema}\\
$G,\ \alpha,\ c$ & Matriz de Gram, columna de productos internos y coeficientes de proyección & \pageref{ampl:gram}\\
$A^*$ & Transpuesta conjugada de una matriz compleja & \pageref{ampl:gram}\\
$b_n^*,\ S_N$ & Funcionales coordenados y sumas parciales de una base de Schauder & \pageref{ampl:schauder}\\
\end{longtable}
\normalsize
